% musixjt.tex :Jazz Text font definitions for MusiXTeX
%
% usage: 
%
%    \input musixtex
%    \input musixjt
%    ...
%
%   MusiXjt.tex is free software; you can redistribute it and/or modify
%   it under the terms of the GNU General Public License as published by
%   the Free Software Foundation; either version 2, or (at your option)
%   any later version.
%
%   MusiXjt.tex is distributed in the hope that it will be useful,
%   but WITHOUT ANY WARRANTY; without even the implied warranty of
%   MERCHANTABILITY or FITNESS FOR A PARTICULAR PURPOSE.  See the
%   GNU General Public License for more details.
%
%   You should have received a copy of the GNU General Public License
%   along with MusiXTeX; see the file COPYING.  If not, write to
%   the Free Software Foundation, Inc., 59 Temple Place - Suite 330,
%   Boston, MA 02111-1307, USA.
%
%   Copyright 2015-2017  Bob Tennent rdt@cs.queensu.ca
%
\ifx\undefined\startpiece\errmessage{Input musixtex.tex before musixjt.tex}\fi
\immediate\write16{MusiXjt\space<2021/08/28>}
%

% 7pt 
\font\sevenrm=MuseJazzText at 7pt
\font\sevenbf=MuseJazzText at 7pt
\font\sevenit=MuseJazzText at 7pt
\font\sevenbi=MuseJazzText at 7pt
\font\sevensc=MuseJazzText at 7pt
% 8pt 
\font\eightrm=MuseJazzText at 8pt
\font\eightbf=MuseJazzText at 8pt
\font\eightit=MuseJazzText at 8pt
\font\eightbi=MuseJazzText at 8pt
\font\eightsc=MuseJazzText at 8pt
% 9pt
\font\ninerm=MuseJazzText at 9pt
\font\ninebf=MuseJazzText at 9pt
\font\nineit=MuseJazzText at 9pt
\font\ninebi=MuseJazzText at 9pt
\font\ninesc=MuseJazzText at 9pt
% 10pt
\font\tenrm=MuseJazzText at 10pt
\font\tenbf=MuseJazzText at 10pt
\font\tenit=MuseJazzText at 10pt
\font\tenbi=MuseJazzText at 10pt
\font\tensc=MuseJazzText at 10pt
% 11pt
\font\elevenrm=MuseJazzText at 11pt
\font\elevenbf=MuseJazzText at 11pt
\font\elevenit=MuseJazzText at 11pt
\font\elevenbi=MuseJazzText at 11pt
\font\elevensc=MuseJazzText at 11pt
% 12pt
\font\twelverm=MuseJazzText at 12pt
\font\twelvebf=MuseJazzText at 12pt
\font\twelveit=MuseJazzText at 12pt
\font\twelvebi=MuseJazzText at 12pt
\font\twelvesc=MuseJazzText at 12pt
% 14pt
\font\frtrm=MuseJazzText at 14pt
\font\frtbf=MuseJazzText at 14pt
\font\frtit=MuseJazzText at 14pt
\font\frtbi=MuseJazzText at 14pt
\font\frtsc=MuseJazzText at 14pt
% 17pt
\font\svtrm=MuseJazzText at 17pt
\font\svtbf=MuseJazzText at 17pt
\font\svtit=MuseJazzText at 17pt
\font\svtbi=MuseJazzText at 17pt
\font\svtsc=MuseJazzText at 17pt
% 20pt
\font\twtyrm=MuseJazzText at 20pt
\font\twtybf=MuseJazzText at 20pt
\font\twtyit=MuseJazzText at 20pt
\font\twtybi=MuseJazzText at 20pt
\font\twtysc=MuseJazzText at 20pt
% 25pt
\font\twfvrm=MuseJazzText at 25pt
\font\twfvbf=MuseJazzText at 25pt
\font\twfvit=MuseJazzText at 25pt
\font\twfvbi=MuseJazzText at 25pt
\font\twfvsc=MuseJazzText at 25pt
%
% large fonts for titles
%
\def\bigtype{\bigfont}
\def\Bigtype{\Bigfont}
\def\BIgtype{\BIgfont}
\def\BIGtype{\BIGfont}

%
\font\ppfftwelve=MuseJazzText at 8pt
\font\ppffsixteen=MuseJazzText at 10pt
\font\ppfftwenty=MuseJazzText at 12pt
\font\ppfftwentyfour=MuseJazzText at 14pt
\font\ppfftwentynine=MuseJazzText at 17pt
%
\def\f{{\ppff f}}
\def\ff{{\ppff ff}}
\def\fp{{\ppff fp}}
\def\fff{{\ppff fff}}
\def\ffff{{\ppff ffff}}
\def\mf{{\ppff mf}}
\def\p{{\ppff p}}
\def\pp{{\ppff pp}}
\def\ppp{{\ppff ppp}}
\def\pppp{{\ppff pppp}}
%
\edef\catcodeat{\the\catcode`\@}\catcode`\@=11
%
\def\sF{{\ppff s\p@kern f}}
\def\sfz{{\ppff s\p@kern f\f@kern z}}
\def\sfzp{{\ppff s\p@kern f\f@kern z\p@kern p}}

\def\mp@{{\ppff mp}}
\let\mezzopiano\mp@

\catcode`\@=\catcodeat

%
% Redefine accented characters for 8-bit font, suggested by David Carlisle:
%

\ifx\documentclass\undefined
\catcode`\@=11
\def\ProvidesFile#1[#2]{}
\def\DeclareFontEncoding#1#2#3{}
\def\DeclareTextAccent#1#2#3{%
\def#1##1{%
\expandafter\ifx\csname T1\string#1-\string##1\endcsname\relax
{\accent#1 ##1}%
\else
\csname T1\string#1-\string##1\expandafter\endcsname
\fi}}
\def\DeclareTextCommand#1#2{\xdtcmd}%not today
\def\xdtcmd#1#{\xxdtcmd}%not today
\def\xxdtcmd#1{}%not today
\def\DeclareTextCompositeCommand#1#2#3#4{}%not today
\def\DeclareTextSymbol#1#2#3{%
\def#1{\char#3\relax}}
\def\DeclareTextComposite#1#2#3#4{%
\expandafter\def\csname T1\string#1-\string#3\endcsname{\char#4\relax}}

\input t1enc.def 

% \c needs special treatment
\def\c#1{\leavevmode\ifx c#1\char231 \else\setbox\z@\hbox{#1}\ifdim\ht\z@=1ex\accent11 #1%
     \else{\ooalign{\unhbox\z@\crcr
        \hidewidth\char11\hidewidth}}\fi\fi}
\catcode`\@=\catcodeat
\fi

%  tt fonts needed by musixsty
%
\font\eighttt=ectt0800
\font\ninett=ectt0900
\font\tentt=ectt1000
\font\twelvett=ectt1200
\font\frttt=ectt1440
\font\svttt=ectt1728
\font\twtytt=ectt2074
\font\twfvtt=ectt2488
%
%  sl fonts needed by musixsty
%
\font\eightsl=ecsl0800
\font\ninesl=ecsl0900
\font\tensl=ecsl1000
\font\twelvesl=ecsl1200
\font\frtsl=ecsl1440
\font\svtsl=ecsl1728
\font\twtysl=ecsl2074
\font\twfvsl=ecsl2488

\normtype
\endinput
